\documentclass[11pt,a4paper]{article}
\usepackage[utf8]{inputenc}
\usepackage{amsmath,amssymb,amsthm}
\usepackage{listings}
\usepackage{geometry}
\usepackage{hyperref}
\geometry{margin=2.5cm}

\newtheorem{theorem}{Theorem}
\newtheorem{lemma}[theorem]{Lemma}

\lstset{
  language=C++,
  basicstyle=\ttfamily\small,
  keywordstyle=\bfseries,
  breaklines=true,
  frame=single,
  numbers=left,
  numberstyle=\tiny,
  tabsize=4,
  showstringspaces=false
}

\title{IOI 2006 -- Beans (Mexicane\~{n}o)}
\author{Solution}
\date{}

\begin{document}
\maketitle

\section{Problem Statement}
We are given $N$ beans with integer values $a_1, a_2, \ldots, a_N$ arranged in a circle. We must select a subset of beans such that no two selected beans are adjacent in the circle, maximizing the total value of selected beans.

\section{Solution}

\subsection{Linear Case (No Wrap-Around)}
Consider first the simpler problem where the beans are arranged in a line. Define:
\begin{align*}
dp[i][0] &= \text{maximum sum using beans } 1, \ldots, i \text{ with bean } i \text{ \emph{not} selected,} \\
dp[i][1] &= \text{maximum sum using beans } 1, \ldots, i \text{ with bean } i \text{ selected.}
\end{align*}
The recurrence is:
\begin{align*}
dp[i][0] &= \max(dp[i-1][0],\; dp[i-1][1]), \\
dp[i][1] &= dp[i-1][0] + a_i,
\end{align*}
with base case $dp[1][0] = 0$, $dp[1][1] = a_1$. The answer for the linear case is $\max(dp[N][0], dp[N][1])$.

This can be computed using only $O(1)$ space since each row depends only on the previous row.

\subsection{Circular Case}
When the beans are arranged in a circle, beans $1$ and $N$ are adjacent. We reduce to two linear sub-problems:
\begin{enumerate}
    \item \textbf{Exclude bean 1}: Solve the linear problem on beans $2, 3, \ldots, N$.
    \item \textbf{Include bean 1}: Since bean 1 is selected, beans 2 and $N$ are forbidden. Solve the linear problem on beans $3, 4, \ldots, N-1$ and add $a_1$.
\end{enumerate}
The answer is the maximum of these two cases.

\begin{lemma}
This reduction is correct because every valid circular selection either includes bean 1 or does not.
\end{lemma}

\subsection{Complexity}
\begin{itemize}
    \item \textbf{Time:} $O(N)$ -- two linear passes.
    \item \textbf{Space:} $O(1)$ with the rolling-variable optimization (or $O(N)$ if storing the full DP table).
\end{itemize}

\section{C++ Solution}

\begin{lstlisting}
#include <bits/stdc++.h>
using namespace std;

int main(){
    ios::sync_with_stdio(false);
    cin.tie(nullptr);

    int N;
    cin >> N;
    vector<long long> a(N);
    for(int i = 0; i < N; i++) cin >> a[i];

    if(N == 1){
        cout << max(0LL, a[0]) << "\n";
        return 0;
    }
    if(N == 2){
        cout << max({0LL, a[0], a[1]}) << "\n";
        return 0;
    }

    // Solve max-weight independent set on a path
    auto solveLinear = [](const long long* v, int n) -> long long {
        if(n <= 0) return 0;
        if(n == 1) return max(0LL, v[0]);
        long long prev_no = 0, prev_yes = max(0LL, v[0]);
        for(int i = 1; i < n; i++){
            long long new_no = max(prev_no, prev_yes);
            long long new_yes = prev_no + v[i];
            prev_no = new_no;
            prev_yes = new_yes;
        }
        return max(prev_no, prev_yes);
    };

    // Case 1: exclude a[0], solve on a[1..N-1]
    long long ans1 = solveLinear(a.data() + 1, N - 1);

    // Case 2: include a[0], exclude a[1] and a[N-1], solve on a[2..N-2]
    long long ans2 = a[0] + solveLinear(a.data() + 2, N - 3);

    cout << max(ans1, ans2) << "\n";
    return 0;
}
\end{lstlisting}

\end{document}
